\chapter{Custom node parser theo cấu trúc văn bản pháp quy}
\label{ch:node-parser}

\section{Vì sao chọn}

Văn bản quy chế chia theo Chương --- Điều --- khoản, và văn bản sửa đổi (QĐ 3183) chia theo ``mục sửa
đổi Điều $N$''. Câu trả lời và trích dẫn đều theo Điều, nên Node phải trùng ranh giới Điều. Bộ cắt mặc
định của LlamaIndex (\texttt{SentenceSplitter}) không biết ranh giới này.

Giai đoạn 1 của bản báo cáo trước dùng \texttt{src/ingestion.py}: OCR từng trang và đóng gói
\emph{mỗi trang thành một \texttt{Document}}. Cách đó có hai nhược điểm đã đo được: (1) một Điều
nằm vắt qua hai trang bị chia đôi; (2) phần ``Quyết định ban hành'' ở trang đầu sinh ra các Node
trùng số Điều với phần Quy chế (Điều 1--3 xuất hiện hai lần), làm hỏng phép chấm truy hồi theo số
Điều. Hai module \texttt{src/ingestion.py} và \texttt{src/node\_parser.py} được giữ lại trong
repository để đối chiếu nhưng \textbf{không còn nằm trong pipeline}; dữ liệu vào Node ngày nay đi
qua tầng 2 nêu trên.

\section{Cơ chế}

Cơ chế tách theo heading được mô tả ở Mục~\ref{sec:node}. Phần dưới là hai quyết định về metadata ---
nơi custom parser khác hẳn parser mặc định.

\paragraph{Loại metadata rỗng.} Bản trước gắn cho \emph{mọi} Node bốn khóa \texttt{van\_ban},
\texttt{chuong}, \texttt{dieu}, \texttt{sua\_doi\_dieu}, kể cả khi giá trị rỗng. Vì LlamaIndex ghép
metadata vào văn bản trước khi embedding, mỗi vector đều chứa các dòng vô nghĩa kiểu
\texttt{chuong:}, \texttt{dieu:} trong khi kho chỉ có vài chục Node. Bản mới lọc bỏ mọi khóa rỗng.

\paragraph{Chỉ một trường được embedding.} Bản mới sinh thêm khóa \texttt{nguon\_trich} và chỉ trường
này được đưa vào embedding; các khóa kỹ thuật còn lại bị loại bằng
\texttt{excluded\_embed\_metadata\_keys}. Văn bản đưa vào bộ mã hóa vì thế có dạng:

\begin{lstlisting}[caption={Văn bản đưa vào embedding của một Node (in bằng \texttt{metadata\_mode="embed"}).}, label={lst:embed-text}]
nguon_trich: Điều 9 — Quy chế đào tạo trình độ đại học SGU 2021

<nguyên văn Điều 9>
\end{lstlisting}

Trước đây cùng Node này cho chuỗi \texttt{van\_ban: ...\textbackslash nchuong: \textbackslash ndieu: 9\textbackslash nsua\_doi\_dieu:} ---
vừa nhiễu vừa chứa khóa tiếng Anh không dấu. Chỉ số \texttt{nguon\_trich} cũng chính là nhãn
\texttt{node\_label()} dùng để in nguồn, nên phần hiển thị và phần embedding không còn lệch nhau.
Đây là ví dụ cho thấy một quyết định nhỏ ở tầng Node ảnh hưởng trực tiếp tới chất lượng vector.

\subsection{Mã nguồn}

\begin{lstlisting}[language=Python, caption={Gắn metadata đã lọc rỗng và chọn trường đưa vào embedding (\texttt{src/processed\_loader.py}).}, label={lst:gan-metadata}]
def _gan_metadata(node: TextNode, meta: dict[str, str | int]) -> TextNode:
    """Gắn metadata đã lọc rỗng + `nguon_trich`, và đặt đúng trường nào được embed / vào prompt."""
    sach = {k: v for k, v in meta.items() if str(v).strip()}
    sach["nguon_trich"] = nguon_trich(sach)
    node.metadata = sach
    node.excluded_embed_metadata_keys = [k for k in sach if k != "nguon_trich"]
    node.excluded_llm_metadata_keys = [k for k in KHONG_LLM if k in sach]
    return node
\end{lstlisting}

\section{Đã làm gì}

\begin{itemize}
  \item Viết \texttt{src/processed\_loader.py} thay cho \texttt{node\_parser.py} (regex \texttt{DIEU\_PATTERN} trên
        văn bản thô) --- nay dựa vào heading do \texttt{clean\_md.py} dựng.
  \item Tách riêng phần ``Quyết định ban hành'' (\texttt{phan="ban\_hanh"}) để Điều 1--3 của quyết định không
        trùng nhãn với Điều 1--3 của quy chế.
  \item Gắn \texttt{LOADER\_VERSION} để chỉ mục tự build lại khi logic cắt đổi (Chương~\ref{ch:persist}).
\end{itemize}

\section{Số liệu}

54 Node từ 6 văn bản (Mục~\ref{sec:ingestion-ket-qua}); chỉ mục và loader khớp 54/54
(Mục~\ref{sec:kiem-chung-moi-truong}). Truy hồi đạt R@1 81\% trên bộ 01, cao hơn BM25 23 điểm phần trăm
(Mục~\ref{sec:danh-gia-thuc-nghiem}).

\chuatrienkhai{Bảng so sánh cách chunk (B1) --- xem Mục~\ref{sec:node}.}

\section{Kiểm thử}

\texttt{tests/test\_processed\_loader.py} khóa lại đúng những hành vi trên: số Node và nhãn Điều,
việc loại metadata rỗng, việc chỉ \texttt{nguon\_trich} được embedding, và giá trị của
\texttt{nguon\_trich} cho ba trường hợp (Điều thường, mục sửa đổi, phần ban hành).
\texttt{tests/test\_corpus.py} kiểm tra trên corpus thật: mọi tiền tố trong danh sách văn bản khớp
đúng một file, và mọi Node đều có \texttt{nguon\_trich}. Chạy \texttt{make test}: 26 kiểm thử, không
cần GPU và không cần Ollama.

\section{Hạn chế đã biết}

\begin{itemize}
  \item Cách cắt phụ thuộc \texttt{clean\_md.py} dựng heading đúng. Nếu một văn bản bị OCR làm mất
        dòng ``Điều $N$.'' thì Điều đó bị gộp vào Điều trước hoặc vào Node Chương.
  \item ``Token'' vẫn đếm bằng số từ tách theo khoảng trắng, chỉ là xấp xỉ.
  \item Hai văn bản một trang cho mỗi văn bản đúng một Node dài, chưa cắt theo khoản.
\end{itemize}

\section{Dùng ở đâu trong đồ án}

Mọi bước sau đều đi trên Node này: chỉ mục, truy hồi, nhãn nguồn \texttt{node\_label()}, và luật chấm
đúng/sai theo (văn bản, Điều) của bộ đánh giá.
